% Chapter — Harris–Priester upper-atmosphere density model (FR-8). Follows
% the mandatory FR-29 six-section template: purpose; assumptions and domain
% of validity; derivation; discretization and implementation; validation
% evidence; references. The equation labels eq:hp:interp, eq:hp:cospow,
% eq:hp:density, eq:hp:apex, and eq:hp:geodetic are echoed verbatim in
% comments in cpp/src/models/atmosphere_hp.cpp (FR-29
% equation-label-to-code traceability); renaming them breaks that trace.
\chapter{Earth Atmosphere: Harris--Priester Density Model}
\label{ch:harrispriester}

\section{Purpose}
\label{sec:hp:purpose}

This chapter documents the orbit-decay drag atmosphere (FR-8): the modified
Harris--Priester density model in the form given by Montenbruck and
Gill~\cite[Sect.~3.5.2]{montenbruck2000}, whose underlying physics is the
time-dependent upper-atmosphere model of Harris and
Priester~\cite{harris1962}. The model is fully deterministic --- a fixed
mean-solar-activity coefficient table plus a diurnal-bulge term, with no
space-weather inputs (D-10 determinism; FR-8 explicitly excludes
space-weather-driven models). It supplies the density consumed by the
cannonball drag acceleration of Chapter~\ref{ch:drag} in the orbital
regime, 100--1000~km. Formulation compatibility with Orekit's
\texttt{HarrisPriester} class~\cite{orekithp} is a design requirement: the
frozen Phase~3 cross-tool drag case (exit criterion~5, decision D-15) uses
Orekit as truth because GMAT does not implement Harris--Priester. The model
is implemented in \texttt{cpp/src/models/atmosphere\_hp.cpp} with its
interface in \texttt{cpp/include/star/models/atmosphere\_hp.hpp}.

\section{Assumptions and domain of validity}
\label{sec:hp:domain}

Assumptions:
\begin{enumerate}
  \item \textbf{Mean solar activity.} The committed coefficient table is
    the Montenbruck--Gill mean-activity table; actual thermospheric
    density varies by factors of several with the solar cycle and
    geomagnetic activity. Montenbruck and Gill place the model's
    representative accuracy near $10$--$20\,\%$ for the regimes it was
    fitted to~\cite[Sect.~3.5.2]{montenbruck2000}; treat absolute decay
    predictions accordingly.
  \item \textbf{Diurnal bulge only.} The single harmonic
    $\cos^{n}(\psi/2)$ apex/antapex modulation is the only local-time
    dependence; semiannual, 27-day, and geomagnetic variations are not
    modeled.
  \item \textbf{Static bulge geometry.} The bulge apex is taken at the
    Sun's declination, lagging the Sun by \ang{30} in right ascension
    (the thermospheric response lag)~\cite[Sect.~3.5.2]{montenbruck2000};
    the same constant is used by the Orekit reference
    implementation~\cite{orekithp}.
  \item \textbf{Geodetic height argument.} The table altitude is geodetic
    height above the reference ellipsoid. This project evaluates it above
    the WGS84 ellipsoid ($a = \qty{6378137.0}{\metre}$,
    $1/f = 298.257223563$~\cite{nima2000}; values in
    \texttt{star/constants.hpp}), matching the Orekit cross-tool
    configuration. The density function itself takes the altitude as an
    argument, so a caller may supply a height computed against any
    ellipsoid.
  \item \textbf{Inclination-dependent exponent.} The bulge exponent $n$
    ranges from 2 for low-inclination orbits to 6 for polar
    orbits~\cite[Sect.~3.5.2]{montenbruck2000}. The default is $n = 4$,
    adopted from the Orekit reference implementation~\cite{orekithp} so
    that default configurations of the two tools coincide; mission
    configuration overrides it per orbit.
\end{enumerate}

Domain of validity: geodetic altitude
$h \in [\qty{100}{\kilo\metre}, \qty{1000}{\kilo\metre}]$ (the table
span) and $n \in [2, 6]$. Out-of-domain behavior, matching the code
exactly: below \qty{100}{\kilo\metre} and for non-finite or
out-of-range arguments ($|\cos\psi| > 1$, $n < 2$, $n > 6$) the model
throws \texttt{std::domain\_error}; above \qty{1000}{\kilo\metre} it
returns exactly zero. Both choices mirror the Orekit reference
implementation (which raises an altitude exception below the table and
returns zero above it), keeping the cross-tool trajectories comparable;
they are gated by \testid{ATM-HP-DOMAIN}.

\section{Derivation}
\label{sec:hp:derivation}

\subsection{Antapex and apex profiles}

The committed table lists, at altitudes $h_i$ (100--1000~km, 50 rows), the
minimum (antapex) density $\rho_{m,i}$ and maximum (apex) density
$\rho_{M,i}$ for mean solar activity
\cite[Sect.~3.5.2, pp.~89--91]{montenbruck2000}. Within a bracket
$[h_i, h_{i+1}]$ each profile is an exponential with the constant scale
height implied by the bracket endpoints,
$H_m = (h_i - h_{i+1}) / \ln(\rho_{m,i+1}/\rho_{m,i})$, so that
\begin{equation}
  \rho_m(h) \;=\; \rho_{m,i} \exp\!\left(\frac{h_i - h}{H_m}\right)
  \;=\; \rho_{m,i}
    \left(\frac{\rho_{m,i+1}}{\rho_{m,i}}\right)^{\textstyle
      \frac{h_i - h}{h_i - h_{i+1}}},
  \label{eq:hp:interp}
\end{equation}
and identically for $\rho_M(h)$. The second (power) form is exactly the
first with the scale height eliminated; it is the form the reference
implementation evaluates~\cite{orekithp}, and this implementation follows
it operation for operation so the two tools differ only in floating-point
rounding. At $h = h_i$ the exponent is zero and the tabulated value is
returned exactly.

\subsection{Diurnal bulge}

Let $\vv{u}_r$ be the unit satellite position vector and $\vv{u}_b$ the
unit vector toward the bulge apex, both resolved in any common
Earth-equatorial frame. With $\cos\psi = \vv{u}_b \cdot \vv{u}_r$, the
half-angle identity gives the bulge weight without evaluating $\psi$
itself:
\begin{equation}
  \cos^{2}\!\frac{\psi}{2} \;=\; \frac{1 + \cos\psi}{2},
  \qquad
  \cos^{n}\!\frac{\psi}{2}
    \;=\; \left(\frac{1 + \cos\psi}{2}\right)
          \left(\sqrt{\frac{1 + \cos\psi}{2}}\,\right)^{\!n-2},
  \label{eq:hp:cospow}
\end{equation}
which is well defined for the non-integer exponents the configuration
permits. The density is the antapex profile plus the bulge modulation
\cite[Sect.~3.5.2]{montenbruck2000}:
\begin{equation}
  \rho(h, \psi) \;=\; \rho_m(h)
    + \bigl[\rho_M(h) - \rho_m(h)\bigr]\cos^{n}\!\frac{\psi}{2}.
  \label{eq:hp:density}
\end{equation}
At $\psi = 0$ (apex) the model returns $\rho_M(h)$; at $\psi = \pi$
(antapex) it returns $\rho_m(h)$ --- the two cases the criterion-8 gate
pins.

The apex direction is the unit Sun direction $\vv{u}_s = (s_x, s_y, s_z)$
rotated by the lag angle $\lambda_b = \ang{30}$ about the Earth's polar
($+z$) axis, which advances its right ascension by $\lambda_b$ while
preserving its declination:
\begin{equation}
  \vv{u}_b \;=\;
  \begin{bmatrix}
    s_x \cos\lambda_b - s_y \sin\lambda_b \\
    s_x \sin\lambda_b + s_y \cos\lambda_b \\
    s_z
  \end{bmatrix}.
  \label{eq:hp:apex}
\end{equation}

\subsection{Geodetic altitude}

The model's height argument is geodetic. For callers holding an
Earth-fixed Cartesian position $(x, y, z)$, the module provides the
closed-form geodetic-altitude approximation of
Bowring~\cite{bowring1976}: with $p = \sqrt{x^2 + y^2}$,
$e^2 = f(2-f)$, $e'^2 = e^2/(1-e^2)$, $b = a(1-f)$, and starting parametric
latitude $\tan u = (z/p)(a/b)$, one fixed refinement of
\begin{equation}
  \tan\phi \;=\; \frac{z + e'^2\, b \sin^3 u}{p - e^2\, a \cos^3 u},
  \qquad
  h \;=\; \frac{p}{\cos\phi} - \frac{a}{\sqrt{1 - e^2 \sin^2\phi}},
  \label{eq:hp:geodetic}
\end{equation}
after which $u$ is recomputed from $\phi$ and
equation~\eqref{eq:hp:geodetic} is applied once more. Bowring's iteration
converges cubically; for heights within this model's domain the fixed
two-pass evaluation is accurate to well below a millimetre, which is seven
orders of magnitude below the density model's own fidelity. The fixed pass
count keeps the evaluation branch-free and deterministic (D-10).

\section{Discretization and implementation}
\label{sec:hp:implementation}

\begin{enumerate}
  \item \textbf{Module and traceability.} The model lives in
    \texttt{cpp/src/models/atmosphere\_hp.cpp}; the source comments echo
    the labels \texttt{eq:hp:interp}, \texttt{eq:hp:cospow},
    \texttt{eq:hp:density}, \texttt{eq:hp:apex}, and
    \texttt{eq:hp:geodetic} verbatim.
  \item \textbf{Compiled-in table with a committed check copy.} The
    50-row coefficient table is compiled into the model source in
    \unit{\kilogram\per\metre\cubed} (the book prints
    \unit{\gram\per\kilo\metre\cubed} $= 10^{-12}$
    \unit{\kilogram\per\metre\cubed}), digit-for-digit identical to the
    Orekit reference table. The committed golden
    \texttt{tests/golden/atmosphere/harris\_priester\_table.toml} is the
    independent transcription check: test \testid{ATM-HP-NODES} requires
    the compiled table to match it bit-for-bit after \texttt{strtod}. A
    runtime file load was rejected: the core never parses text (D-2).
  \item \textbf{Purity.} The density function is a pure function of
    (geodetic altitude, $\cos\psi$, $n$); the bulge-apex construction
    (equation~\eqref{eq:hp:apex}) and the geodetic-altitude helper
    (equation~\eqref{eq:hp:geodetic}) are separate pure functions. Frame
    handling and ellipsoid choice remain with the caller (the EOM layer
    owns frames), which is what keeps the model testable against the
    published values without an ephemeris in the loop.
  \item \textbf{Evaluation-order compatibility.} Bracket search is a
    fixed-order forward scan; the interpolation is evaluated in the power
    form of equation~\eqref{eq:hp:interp}; the bulge weight uses
    equation~\eqref{eq:hp:cospow} with the reference implementation's
    guard (weight set to zero when $\cos(\psi/2) \le 10^{-12}$, i.e.\
    within double rounding of the exact antapex) --- each choice mirrors
    Orekit~\cite{orekithp} so cross-tool differences reduce to
    floating-point rounding.
  \item \textbf{Determinism.} Straight-line binary64 arithmetic, fixed
    scan order, no data-dependent iteration counts; libm
    \texttt{pow}/\texttt{log}/\texttt{sqrt} contribute the usual
    last-ulp cross-platform spread, orders below every committed
    tolerance.
\end{enumerate}

\section{Validation evidence}
\label{sec:hp:validation}

Table~\ref{tab:hp:validation} lists the tests that constitute this
chapter's validation evidence. Each test identifier is machine-checked to
exist in the test tree by \texttt{scripts/lint\_chapters.py}; golden
provenance and tolerances are recorded in
\texttt{tests/golden/atmosphere/manifest.toml}.

\begin{table}[htbp]
  \centering
  \caption{Validation evidence for the Harris--Priester model (doctest,
    \texttt{cpp/tests/test\_atmosphere.cpp}).}
  \label{tab:hp:validation}
  \begin{tabular}{lp{8.6cm}}
    \toprule
    Test ID & Acceptance basis \\
    \midrule
    \testid{ATM-HP-NODES} &
      Phase~3 exit criterion~8: with the bulge term pinned to its minimum
      ($\cos\psi = -1$) and separately its maximum ($\cos\psi = +1$), the
      model evaluated at every one of the 50 published altitude nodes
      reproduces the published $\rho_{\min}$/$\rho_{\max}$ value to 4
      significant figures (at-node evaluation is exact by construction,
      equation~\eqref{eq:hp:interp}); the compiled table equals the
      committed transcription check copy bit-for-bit. \\
    \testid{ATM-HP-OFFNODE} &
      Eight off-node cases (101--995~km, $\cos\psi \in [-1,1]$,
      $n \in \{2,4,6\}$) match the independent 50-digit mpmath reference
      to $10^{-13}$ relative; the apex construction
      (equation~\eqref{eq:hp:apex}) and the geodetic helper
      (equation~\eqref{eq:hp:geodetic}, closure against exact synthesized
      geodetic points) are exercised in the same case. \\
    \testid{ATM-HP-DOMAIN} &
      Out-of-domain behavior of Section~\ref{sec:hp:domain}:
      \texttt{std::domain\_error} below \qty{100}{\kilo\metre} and for
      invalid $\cos\psi$/$n$/non-finite inputs; exactly zero above
      \qty{1000}{\kilo\metre}. \\
    \bottomrule
  \end{tabular}
\end{table}

\section{References}
\label{sec:hp:references}

This chapter relies on Montenbruck and Gill~\cite{montenbruck2000},
Sect.~3.5.2, for the model formulation and the mean-solar-activity
coefficient table; on Harris and Priester~\cite{harris1962} for the
underlying physical model; on the Orekit \texttt{HarrisPriester}
implementation~\cite{orekithp} as the formulation-compatibility target and
table transcription source (decision D-15); on NIMA
TR8350.2~\cite{nima2000} for the WGS84 ellipsoid parameters; and on
Bowring~\cite{bowring1976} for the geodetic-altitude conversion. Full
bibliographic details appear in the bibliography.
